% ============================================================
%  PREÂMBULO — Dev's Café Biblioteca Didática
% ============================================================

% === CODIFICAÇÃO E LÍNGUA ===
\usepackage[utf8]{inputenc}     % Codificação UTF-8 (pdfLaTeX)
\usepackage[T1]{fontenc}        % Codificação de saída T1
\usepackage[brazilian]{babel}   % Idioma: Português do Brasil

% === TIPOGRAFIA ===
\usepackage{lmodern}
\usepackage{microtype}

% === LAYOUT ===
\usepackage[
  a4paper,
  top=3cm, bottom=3cm,
  left=3cm, right=2.5cm
]{geometry}
\usepackage{setspace}
\onehalfspacing

% === MATEMÁTICA ===
\usepackage{amsmath}
\usepackage{amssymb}
\usepackage{amsthm}
\usepackage{bm}                 % Para negrito em matemática (vetores)

% === CÓDIGO-FONTE ===
\usepackage{listings}
\usepackage{xcolor}

% === GRÁFICOS E FIGURAS ===
\usepackage{graphicx}
\usepackage{float}
\usepackage{caption}
\usepackage{subcaption}
\usepackage{pagecolor}
\usepackage{tikz}               % Para diagramas técnicos

% === TABELAS ===
\usepackage{booktabs}
\usepackage{tabularx}

% === CAIXAS DIDÁTICAS ===
\usepackage{tcolorbox}
\tcbuselibrary{most}

% === ÍCONES ===
\usepackage{fontawesome5}

% === CABEÇALHOS E RODAPÉS ===
\usepackage{fancyhdr}
\pagestyle{fancy}
\fancyhf{}
\fancyhead[L]{\small\nouppercase{\leftmark}}
\fancyhead[R]{\small Dev's Café}
\fancyfoot[C]{\thepage}
\renewcommand{\headrulewidth}{0.4pt}

% === PALETA DE CORES (Dev's Café) ===
\definecolor{devcafe-blue}{HTML}{2563EB}
\definecolor{devcafe-green}{HTML}{16A34A}
\definecolor{devcafe-red}{HTML}{DC2626}
\definecolor{devcafe-yellow}{HTML}{CA8A04}
\definecolor{devcafe-gray}{HTML}{64748B}
\definecolor{devcafe-light}{HTML}{F1F5F9}
\definecolor{devcafe-dark}{HTML}{1E293B}

% === HIPERLINKS ===
\usepackage[
  colorlinks=true,
  linkcolor=devcafe-blue,
  urlcolor=devcafe-blue,
  citecolor=devcafe-gray
]{hyperref}

% === BIBLIOGRAFIA (biblatex + biber) ===
\usepackage[
  backend=biber,
  style=authoryear,
  sorting=nyt,
  maxbibnames=3,
  minbibnames=1,
  maxcitenames=2,
  mincitenames=1,
  giveninits=true,
  uniquename=init,
  dashed=false,
  doi=true,
  isbn=false,
  url=true,
  urldate=comp,
  language=portuguese,
]{biblatex}
\addbibresource{bibliografia.bib}

% === AMBIENTES MATEMÁTICOS ===
\newtheorem{teorema}{Teorema}[chapter]
\newtheorem{lema}[teorema]{Lema}
\newtheorem{proposicao}[teorema]{Proposição}
\newtheorem{corolario}[teorema]{Corolário}
\newtheorem{definicao}{Definição}[chapter]

% === CAIXAS DIDÁTICAS PADRONIZADAS ===

% Dica
\newtcolorbox{dica}{
  colback=devcafe-green!5, colframe=devcafe-green,
  fonttitle=\bfseries,
  title={\faLightbulb\ Dica}
}

% Aviso
\newtcolorbox{aviso}{
  colback=devcafe-yellow!10, colframe=devcafe-yellow,
  fonttitle=\bfseries,
  title={\faExclamationTriangle\ Aviso}
}

% Exemplo
\newtcolorbox{exemplo}[1][Exemplo]{
  colback=white, colframe=devcafe-green,
  fonttitle=\bfseries,
  title={\faCode\ #1}
}

% Exercício
\newtcolorbox{exercicio}[1]{
  colback=devcafe-light, colframe=devcafe-gray,
  fonttitle=\bfseries,
  title={\faEdit\ Exercício: #1}
}

% Fórmula (Nova)
\newtcolorbox{formula}[1]{
  colback=devcafe-blue!5, colframe=devcafe-blue,
  fonttitle=\bfseries,
  title={\faCalculator\ Fórmula: #1}
}

% === CONFIGURAÇÃO DE LISTINGS (GLSL e Pseudocódigo) ===

\lstdefinelanguage{GLSL}{
  morekeywords={attribute, uniform, varying, layout, in, out, inout, float, int, void, bool, mat2, mat3, mat4, vec2, vec3, vec4, ivec2, ivec3, ivec4, bvec2, bvec3, bvec4, sampler1D, sampler2D, sampler3D, samplerCube, struct, for, while, do, if, else, break, continue, return, discard, true, false, precision, highp, mediump, lowp, invariant, smooth, flat, noperspective, texture, texture2D, pow, max, dot, normalize, reflect, refract, clamp},
  sensitive=true,
  morecomment=[l]{//},
  morecomment=[s]{/*}{*/},
  morestring=[b]",
}

\lstset{
  basicstyle=\ttfamily\small,
  keywordstyle=\color{devcafe-blue}\bfseries,
  stringstyle=\color{devcafe-green},
  commentstyle=\color{devcafe-gray}\itshape,
  numberstyle=\tiny\color{devcafe-gray},
  numbers=left,
  stepnumber=1,
  frame=single,
  backgroundcolor=\color{devcafe-light},
  breaklines=true,
  breakatwhitespace=true,
  tabsize=2,
  showstringspaces=false,
  captionpos=b,
  rulecolor=\color{devcafe-gray},
}

\lstdefinelanguage{pseudocode}{
  morekeywords={if, then, else, for, each, in, while, return, function, algorithm, end, and, or, not},
  sensitive=false,
  morecomment=[l]{//},
  morecomment=[s]{/*}{*/},
  morestring=[b]",
}
